\section{Primitive necklaces and divisor-sensitive inversion}

Let $C_n$ be the number of primitive cyclic necklaces in $\Sigma_A$.  Since
$N_n$ counts periodic points rather than necklaces, ordinary M\"obius
inversion gives
\begin{equation}\label{eq:primitive-all}
C_n=\frac1n\sum_{d\mid n}\mu(n/d)N_d.
\end{equation}

For a divisor $d$, set
\[
f(d)=F_{(d+3)/2}\quad(d\text{ odd}),
\]
and define
\[
f_{\EE}(d)=
\begin{cases}f(d),&d\text{ odd},\\G_{d/2},&d\text{ even},\end{cases}
\qquad
f_{\VV}(d)=
\begin{cases}f(d),&d\text{ odd},\\H_{d/2},&d\text{ even}.
\end{cases}
\]

\begin{theorem}[Primitive reflection inversion]\label{thm:primitive}
If $n$ is odd, the number of primitive reversible necklaces is
\begin{equation}\label{eq:primitive-odd}
R_n=\sum_{d\mid n}\mu(n/d)F_{(d+3)/2}.
\end{equation}
If $n$ is even, the two physical axis counts are
\begin{equation}\label{eq:primitive-even}
R_n^{\EE}=\frac12\sum_{d\mid n}\mu(n/d)f_{\EE}(d),
\qquad
R_n^{\VV}=\frac12\sum_{d\mid n}\mu(n/d)f_{\VV}(d),
\end{equation}
and $R_n=R_n^{\EE}+R_n^{\VV}$.
\end{theorem}

\begin{proof}
First subtract all words of proper period by M\"obius inversion while keeping
the reflection operator fixed.  A primitive necklace cannot be stabilized by
two distinct reflections: their product would be a nontrivial stabilizing
rotation.  At odd period, multiplication by two is invertible modulo $n$, so
the $n$ rotated representatives distribute one-to-one across the $n$
reflection axes.  Thus the primitive fixed-word count for one axis already
equals the number of reversible necklaces, proving \eqref{eq:primitive-odd}.

At even period, rotation preserves axis parity.  There are $n/2$ axes of each
parity, and a primitive reversible necklace of a fixed parity contributes two
rotated representatives to every compatible fixed-axis census.  Dividing
the M\"obius-inverted fixed-word count by two proves
\eqref{eq:primitive-even}.  The definition of $f_{\EE},f_{\VV}$ is essential:
an even-period word inherited from an odd proper divisor belongs to the
single odd-period reflection class before it is lifted.
\end{proof}

This proof also explains why a parity-blind necklace formula is wrong.  The
factor $1/2$ and the odd-divisor branch are geometric multiplicities, not
normalization choices.
